% Copyright 2026 The terCAD team. All rights reserved.
% Use of this content is governed by a CC BY-NC-ND 4.0 license that can be found in the LICENSE file.

\subsection{Types of Transactions}
\markboth{Processing}{Types of Transactions}

Accounts and budget categories describe the structure of finances, while transactions bring
them to life. Every financial transaction is recorded as one of three types:

\begin{center}
\small
\begin{tabularx}{\textwidth}{|l|X|X|}
  \hline
  \textbf{Type} & \textbf{Meaning} & \textbf{Effect} \\
  \hline
  Bill & An expense: money leaves an account and is attributed to a budget category. & The
  account decreases; the spent amount of the category increases. \\
  \hline
  Income & Money received on an account: salary, interest, gifts, refunds. & The account
  increases. \\
  \hline
  Transfer & Money moved between own accounts: a cash withdrawal, a credit card repayment, a
  deposit. & One account decreases, another one increases. \\
  \hline
\end{tabularx}
\end{center}

\noindent This distinction is important for the reports. A transfer does not change the net
worth or the result of the month: moving £200 from a debit card to a wallet does not
constitute an expense, and moving it back does not constitute an income. Recording transfers
as bills and incomes would inflate both sides of the summary table and distort the budget
statistics (see \cref{summary:table}).

All three types are created from similar forms that can be accessed via the
\textbf{Create new Transaction} action. This can be opened using the
\textbf{Add Bill, Income or Transfer} button on the main page and on lists of transactions,
or by using the shortcut \key{Ctrl} + \key{N}  on devices with a keyboard.
The form has three tabs -- \textbf{Expense}, \textbf{Income}, and \textbf{Transfer}  -- and
remembers the last selected account, budget category and currency. This means that repeated
entries only require a few taps.


\subsection{Recording a Bill}
\markboth{Processing}{Recording a Bill}

The \textbf{Expense} tab contains the following fields:

\begin{description}
  \item[Account] (required) -- where the money is taken from.
  \item[Currency] and \textbf{Amount of Expenses} (required) -- the amount as it appears on
  the receipt. The currency is preset from the account (or the budget category), but can be
  changed, e.g. for a purchase abroad.
  \item[Description] -- what was bought. A good description helps both the search and the
  prediction of a category.
  \item[Budget Category] (required) -- what the money was spent on.
  \item[Billed At] -- the date and time of the purchase; the current moment by default.
\end{description}


\subsubsection{Category Suggestions}

Below the \textbf{Budget Category} field, the form offers buttons with the most probable
categories for the entered bill. These suggestions are learned from previously recorded bills
on the device, including words from the description, the selected account and the amount
range. Selecting a suggestion with a single tap will apply the category. The more bills that
are recorded, the more precise the suggestions will become.


\subsubsection{Bills in Another Currency}

If the currency of a bill differs from that of the account, budget category or default
currency, an additional block is shown on the form for each such pair.

\begin{description}
  \item[Conversion] -- the exchange rate between the two currencies;
  \item[Taken in \dots] -- the resulting amount in the other currency.
\end{description}

\noindent Both fields are linked, so changing one recalculates the other and vice versa. This
enables you to enter the exact amount charged by the bank (including fees), ensuring that
the account balance matches the statement. The rate is preset using known exchange rates and
can be saved for later use.

For example, if you pay for a dinner costing 45.00 EUR with a card in USD at a rate of
1.0870, 48.92 USD will be taken from your account. If the statement shows 49.30 USD, enter
this amount into the Taken in USD field and the rate will be adjusted automatically.


\subsubsection{Splitting a Bill}

One receipt often covers several categories, such as groceries, household goods and gifts
bought in one supermarket. Rather than recording three bills, record one for the total amount
and split it afterwards (\cref{img:split-bill}).

\begin{enumerate}
  \item open the bill and select \textbf{Split the Bill};
  \item enter the amount of the part, its description, and its budget category;
  \item save -- the part becomes a separate bill, and the amount of the original bill is
  reduced accordingly.
\end{enumerate}

\noindent The parts retain the account and date of the original bill, and their titles are
prefixed with a description of the original (e.g. 'Supermarket: Gift'), making the
relationship visible in any list. Repeat this process for each additional part; anything that
remains will stay in the original category. The total charged to the account does not change.


\subsection{Recording an Income}
\markboth{Processing}{Recording an Income}

The \textbf{Income} tab is a shorter version of the expense form:

\begin{description}
  \item[Account] (required) -- where the money is received.
  \item[Currency] and \textbf{Amount of Income} (required).
  \item[Description] -- the source of the income: salary, freelance, interest, refund.
  \item[Date] -- when the money was received.
\end{description}

\noindent As with bills, an income in a currency other than the currency of the account shows
the \textbf{Conversion} block to define the exact amount credited.

Incomes have no budget category, but they influence budgets with \emph{relative} limits (see
\cref{budget:limit}): each new income within the current interval increases the amount
available in such categories.


\subsubsection{Refunds}

A returned purchase can be recorded in two ways:

\begin{itemize}
  \item \textbf{edit or delete the original bill} -- if the refund arrived in the same period
  and the purchase should disappear from the statistics, as if it had never happened;
  \item \textbf{record an income} -- if the original bill belongs to a closed period, or a
  partial refund (cashback) should be visible as a separate operation.
\end{itemize}


\subsection{Recording a Transfer}
\markboth{Processing}{Recording a Transfer}

The \textbf{Transfer} tab moves money between two own accounts:

\begin{description}
  \item[From Account] and \textbf{To Account} (required) -- the source and the destination.
  \item[Currency] and \textbf{Amount of Transfer} (required).
  \item[Description] -- an optional note.
  \item[Date] -- when the transfer took place.
\end{description}

\noindent Transfers between accounts in different currencies are commonplace, such as buying
foreign currency or paying a credit card bill in a different currency. The Conversion block
determines the exchange rate and the amount credited to the destination account, ensuring
that both balances remain accurate.

\newpage
\noindent Typical transfers:

\begin{itemize}
  \item withdrawing cash: \q{Bank/Debit} $\rightarrow$ \q{Cash};
  \item paying off a credit card: \q{Bank/Debit} $\rightarrow$ \q{Bank/Credit Card};
  \item saving: \q{Bank/Debit} $\rightarrow$ \q{Deposit};
  \item exchanging currency: \q{Cash/EUR} $\rightarrow$ \q{Cash/USD}.
\end{itemize}

\noindent A loan payment usually consists of two parts. The repayment of the principal is
transferred to the 'Credit' account, which reduces the debt. The interest is a bill
representing the actual cost of borrowing.


\subsection{Browsing Transactions}
\markboth{Processing}{Browsing Transactions}

\subsubsection{Lists}

Each type has its own page:

\begin{description}
  \item[Bills] -- the header shows the total expenses of the current month (following the
  \textbf{Start of Month} setting), and the list below is grouped by days, with the date of
  each day pinned on top while scrolling.
  \item[Invoices] -- all incomes.
  \item[Transfers] -- all movements between accounts.
\end{description}

\noindent The lists load progressively while you scroll, so even years of history open
instantly. Selecting a row displays its details. On touch screens, swiping a row provides
quick access to editing and deleting options (\cref{img:bill-view}).

\img{img/ch03-s01-bill__view}{Bills Overview}{img:bill-view}

Transactions can also be accessed from related entities: the details of an account list its
bills, income and transfers, while the details of a budget category list its bills. Selecting
a month in the \textbf{Summary} tab of a category (or in the categories table on the
\textbf{Metrics} page, \cref{img:bill-metrics}) opens all the bills for that category for
that month.

\img{img/ch03-s01-bill__metrics}{Metrics Overview}{img:bill-metrics}


\subsubsection{Search}

The \textbf{Search} action on the pages of bills, invoices, and transfers opens a filter with
the following criteria:

\begin{itemize}
  \item \textbf{Description} -- a part of the text;
  \item \textbf{Account} and \textbf{Account Type} -- e.g. all payments by credit cards;
  \item \textbf{Budget Category} -- for bills only;
  \item \textbf{Currency};
  \item \textbf{From Amount} and \textbf{To Amount} -- a range of amounts;
  \item \textbf{Date Range}.
\end{itemize}

\noindent The criteria can be combined. For example, you could search for all bills over £100
in the 'Leisure' category that were paid by credit card last summer.


\subsubsection{Editing and Deleting}

Any field in a transaction can be edited later via the \textbf{Edit} button. The balances of
the affected accounts and the amounts spent in the affected categories are automatically
recalculated -- including when the account, category, amount or date has been changed.
Internally, changes are stored as new log entries rather than overwrites, which keeps the
history consistent.

\textbf{Delete}: (confirmed by a dialogue box) removes a transaction that has been entered
incorrectly and reverses its effect on balances and budgets. Bear in mind that the balance
date of accounts is updated (see \cref{account:balance}): transactions dated before this date
do not affect the balance, either when they are created or edited or deleted.


\subsection{Recurring Payments}
\markboth{Processing}{Recurring Payments}

Rent, subscriptions, salary, a monthly transfer to savings -- many transactions repeat on a
schedule. Instead of entering them every time, define them once on the \textbf{Automation}
page, in the \textbf{Recurring Payments} tab (see \cref{img:recurring-payments}):

\begin{description}
  \item[Type] -- a bill, an income, or a transfer;
  \item[Payment Interval] -- weekly, monthly, yearly, or
  after some days (for any custom period);
  \item The usual fields of the selected type: account(s), amount, currency, description,
  and budget category.
\end{description}

\noindent When the application is opened, any payments that have become due since the last
launch are created automatically as ordinary transactions, with their scheduled dates. These
can be edited or deleted as required, for example if the electricity bill is higher than
usual. On Android, recurring payments can also be viewed via a widget on the home screen.


\subsection{Importing Transactions}
\markboth{Processing}{Importing Transactions}

There is no need to enter bills and incomes individually. Most banks and financial
applications export statements in one of the formats supported by Fingrom.

\begin{itemize}
  \item \q{CSV} -- a plain table, where the columns are mapped to the fields of a transaction
  during the import;
  \item \q{QIF} -- Quicken Interchange Format;
  \item \q{OFX} -- Open Financial Exchange.
\end{itemize}

\noindent The import is available on the \textbf{Import / Export} tab of the
\textbf{Settings} page. Combined with the \textbf{Balance Date Update} of accounts
(see \cref{account:balance}), it allows bringing years of history without breaking the
current balances.

The procedure is the same for all formats:

\begin{enumerate}
  \item press \textbf{Pick .csv-file}, \textbf{Pick .qif-file}, or \textbf{Pick .ofx-file}
  and select the file;
  \item check the mapping of the columns and the \textbf{Date Format};
  \item define the defaults: the account, the budget category, the currency, and the type for
  the lines that do not specify them;
  \item press \textbf{Parse File} and wait for the confirmation that the process is finished.
\end{enumerate}

\noindent Only bills and incomes are imported; transfers between your own accounts must be
recorded manually. Any accounts or budget categories mentioned in the file that do not yet
exist are created automatically -- review these after the import and adjust their types,
currencies and limits.

\subsubsection{CSV}

A \q{CSV} (comma-separated values) file is a plain table. The first line contains the names
of the columns and each subsequent line represents a transaction. Most banks and spreadsheet
applications can export in this format.

\begin{lstlisting}[caption=Sample CSV File, label=code:import-csv]
Date,Account,Category,Description,Amount,Currency,Type
10/1/2026 09:15,Bank/Debit,Food/Groceries,Supermarket,54.20,EUR,Bill
10/1/2026 18:40,Cash,Food/Dining Out,Coffee,3.50,EUR,Bill
10/5/2026 08:00,Bank/Debit,,Salary,2000.00,EUR,Invoice
10/7/2026 12:30,Bank/Credit Card,Leisure/Hobbies,Books,24.90,EUR,Bill
\end{lstlisting}

\noindent As the names and order of columns differ between banks, each one must be mapped to
a transaction field (the \textbf{Ignore} option skips a column):

\begin{center}
\small
\begin{tabularx}{\textwidth}{|l|X|}
  \hline
  \textbf{Column} & \textbf{Mapping} \\
  \hline
  Date & Bill: Billed At \\
  Account & Account: Title \\
  Category & Budget Category: Title \\
  Description & Bill: Description \\
  Amount & Bill: Amount of Expenses \\
  Currency & Bill: Currency \\
  Type & Bill: Amount Type (Bill, Invoice) \\
  \hline
\end{tabularx}
\end{center}

\noindent A few rules to keep in mind:

\begin{itemize}
  \item \textbf{Type} -- a line with the value \q{Invoice} (as it is named in the language of
  the interface) is imported as an income, any other value as a bill. If the file has no such
  column, the default type is applied to all lines.
  \item \textbf{Amount} -- a positive number with a dot as the decimal separator.
  \item \textbf{Date Format} -- detected automatically from the values of the date column
  (\q{M/d/yyyy HH:mm} in the sample above), and can be corrected manually, e.g.
  \q{dd.MM.yyyy} for \q{31.10.2026}.
  \item Empty cells (like the category of the salary above) are filled with the defaults.
\end{itemize}

\subsubsection{QIF}

\q{QIF} (Quicken Interchange Format) is a text format, in which every line starts with a
one-letter code, and every transaction ends with the \q{\^{}} symbol:

\begin{lstlisting}[caption=Sample QIF File, label=code:import-qif]
!Type:Bank
D10/1/2026
T-54.20
PSupermarket
LFood/Groceries
N1001
^
D10/5/2026
T2000.00
PSalary
N1002
^
\end{lstlisting}

\noindent The codes are mapped automatically:

\begin{center}
\small
\begin{tabularx}{\textwidth}{|l|X|}
  \hline
  \textbf{Code} & \textbf{Meaning} \\
  \hline
  \q{D} & Date of the transaction \\
  \q{T} & Total amount: negative for a bill, positive for an income \\
  \q{P} & Payee, imported as the description \\
  \q{L} & Category, imported as the budget category \\
  \q{N} & Number, imported as the unique identifier of the transaction \\
  \hline
\end{tabularx}
\end{center}

\noindent Other codes (e.g. \q{!Type} in the header, or \q{M} for a memo) are ignored. Since
\q{QIF} does not specify an account, all transactions are imported into the default account.

\subsubsection{OFX}

\q{OFX} (Open Financial Exchange) is a format used by many banks for downloadable statements.
Every transaction is enclosed in a \q{<STMTTRN>} block:

\begin{lstlisting}[language=xml, caption=Sample OFX File, label=code:import-ofx]
OFXHEADER:100
DATA:OFXSGML
VERSION:102

<OFX>
<BANKMSGSRSV1><STMTTRNRS><STMTRS>
<CURDEF>EUR
<BANKTRANLIST>
<STMTTRN>
<TRNTYPE>DEBIT
<DTPOSTED>20261001091500
<TRNAMT>-54.20
<FITID>2026100101
<NAME>Supermarket
</STMTTRN>
<STMTTRN>
<TRNTYPE>CREDIT
<DTPOSTED>20261005080000
<TRNAMT>2000.00
<FITID>2026100501
<NAME>Salary
</STMTTRN>
</BANKTRANLIST>
</STMTRS></STMTTRNRS></BANKMSGSRSV1>
</OFX>
\end{lstlisting}

\noindent The following elements are read from each block:

\begin{center}
\small
\begin{tabularx}{\textwidth}{|l|X|}
  \hline
  \textbf{Element} & \textbf{Meaning} \\
  \hline
  \q{DTPOSTED} & Date and time in the form \q{YYYYMMDDhhmmss} \\
  \q{TRNAMT} & Amount: negative for a bill, positive for an income \\
  \q{NAME} & Description \\
  \q{FITID} & Unique identifier of the transaction \\
  \hline
\end{tabularx}
\end{center}

\noindent Both variants of the format are supported: the older, SGML-based format, where
closing tags of elements are optional, and the newer, XML-based format. As with QIF files,
OFX files do not carry accounts or categories in a form usable by Fingrom, so the defaults
are applied: the account, budget category and currency selected before parsing.


\subsection{Practical Tips}
\markboth{Processing}{Practical Tips}

\textbf{Record right away.} A bill entered at the cash desk takes seconds; a bill
reconstructed at the end of the week is often forgotten or approximated.

\textbf{Write descriptions consistently.} \q{Supermarket}, not \q{supermarket}, \q{SM}, or
\q{groceries at the corner} -- consistent descriptions make search and category suggestions
far more accurate.

\textbf{Use transfers for own money.} Cash withdrawals, card repayments, and savings are
transfers, never bills and incomes.

\textbf{Enter the amount actually charged.} For purchases in a foreign currency, adjust the
\textbf{Taken in \dots} amount to the bank statement, including fees.

\textbf{Split instead of approximating.} One receipt with several categories is one bill
split into parts, not a single bill in the \q{most important} category.

\textbf{Automate the predictable.} Every regular payment defined as recurring is one less
thing to remember.
